% Panorama de MCA_SLOT EXPANDIDO, a nivel de operación, en el estilo de
% panorama_mca_window_expandido: las compuertas y el estado son compuertas y
% comparadores; cada acumulador, un sumador, su mux de carga y su registro con
% la habilitación escrita arriba; la ventana de amplitud, las dos sombras con
% su signo y el pico que quedaría.
% Las redes con nombre (abre, acumula, cierre, abierta, len, peak...) se
% conectan por nombre: el trazo corto que entra a un pin lleva el nombre de la
% red. Las que terminan en _i son puertos de entrada.
% La ventana del pico se dibuja como max(peak, x_c) contra [min, max]; el RTL
% lo escribe como cuatro comparaciones en paralelo, que es lo mismo.
% Fuente: modulos/mca/rtl/mca_slot.sv, instanciado por mca_window como en
% rtl/mine/tops/mca_red_pitaya_top.sv (3 formas del banco).
% Estilo: docs/figuras/estilo_tikz.tex; los colores por función, los mismos de
% panorama_mca_slot.tex. El borde grueso es el límite del módulo.
% El epígrafe no va en la figura: está en panorama_mca_slot_expandido.txt.
%
% Compilar:   pdflatex panorama_mca_slot_expandido.tex
%             pdftocairo -svg panorama_mca_slot_expandido.pdf panorama_mca_slot_expandido.svg
%
\documentclass[border=0pt,tikz]{standalone}
\usepackage[utf8]{inputenc}
\usepackage[T1]{fontenc}
\usepackage{lmodern}
\usepackage[none]{hyphenat}
% Estilo común de las figuras TikZ (panoramas y diagramas de bloques).
%
% Cada figura lo incluye con % Estilo común de las figuras TikZ (panoramas y diagramas de bloques).
%
% Cada figura lo incluye con % Estilo común de las figuras TikZ (panoramas y diagramas de bloques).
%
% Cada figura lo incluye con \input{../../figuras/estilo_tikz.tex} (las figuras
% viven en docs/<módulo>/figuras/). Un solo lugar para la paleta, los estilos de
% bloque y el formato de los parámetros: si cambia algo acá, cambia en todas.
%
% La paleta es la misma de docs/figuras/estilo_svg.py y de PALETA en
% modulos/mk/datasheet.py; si se toca un color, se toca en los tres.
%
% Requiere: \usepackage{tikz} (lo da la clase standalone con la opción tikz).

\usetikzlibrary{arrows.meta,positioning,calc,fit,backgrounds,shapes.multipart,
                decorations.markings,shadows}

% ---------------------------------------------------------------- colores ---
\definecolor{azul}{RGB}{108,172,228}      % acondicionamiento
\definecolor{verde}{RGB}{110,214,142}     % captura / histogramas en BRAM
\definecolor{amarillo}{RGB}{252,232,160}  % trigger
\definecolor{naranja}{RGB}{252,206,160}   % camino a DDR
\definecolor{grisr}{RGB}{226,226,226}     % compartido / infraestructura
\definecolor{rojo}{RGB}{240,128,128}      % bus del sistema

% ---------------------------------------------------------------- estilos ---
\tikzset{
  >={Stealth[length=2.4mm,width=1.8mm]},
  % fondo blanco siempre: el SVG se ve igual en modo claro y oscuro
  fondo/.style   = {show background rectangle, inner frame sep=8pt,
                    background rectangle/.style={fill=white, draw=none}},
  bloque/.style  = {draw=black, line width=0.6pt, rounded corners=2pt,
                    align=center, text width=2.2cm, minimum height=1.15cm,
                    inner sep=2pt, font=\scriptsize, fill=#1},
  datos/.style   = {bloque=azul!35},
  trig/.style    = {bloque=amarillo},
  captura/.style = {bloque=verde!55},
  ddr/.style     = {bloque=naranja!70},
  comp/.style    = {bloque=grisr},
  bus/.style     = {bloque=rojo!45},
  % lógica escrita dentro del módulo, sin submódulo propio
  etapa/.style   = {bloque=#1, line width=0.5pt},
  etapa/.default = verde!20,
  % lo que el bitstream no trae (atado a cero, stub, salida muda)
  mudo/.style    = {bloque=white, draw=black!45, dashed, text=black!55},
  %
  % BLOQUE CON PARÁMETROS, en el formato del datasheet: el encabezado lleva el
  % nombre y los parámetros con que lo instancia el padre (NOMBRE=valor, en
  % monoespaciada, uno por línea); el cuerpo, más claro, lo que hace.
  %   \node[datos, param={azul!35}{azul!12}] {\blk{nombre}{} \prm{A=1\\B=2}
  %                                            \nodepart{two} función};
  param/.style 2 args = {rectangle split, rectangle split parts=2,
                         rectangle split part fill={#1, #2}},
  %
  flujo/.style   = {->, line width=0.8pt},
  ctrl/.style    = {->, line width=0.6pt, black!65},
  % la línea que pasa por encima en un cruce: un halo blanco la separa
  cruce/.style   = {preaction={draw=white, line width=3.4pt, -}},
  % marca de ancho de bus: barra oblicua a mitad de la flecha, con el número
  % encima. \draw[flujo, ancho=14] ...; para otra posición, ancho en=0.3.
  ancho/.style   = {decoration={markings, mark=at position \anchopos with {
                      \draw[line width=0.6pt, -] (-0.1,-0.14) -- (0.1,0.14);
                      \node[font=\scriptsize, above=1.5pt] {#1};}},
                    postaction={decorate}},
  % la misma marca en una flecha vertical: el número va a la derecha
  anchov/.style  = {decoration={markings, mark=at position \anchopos with {
                      \draw[line width=0.6pt, -] (-0.14,-0.1) -- (0.14,0.1);
                      \node[font=\scriptsize, right=3pt] {#1};}},
                    postaction={decorate}},
  ancho en/.code = {\def\anchopos{#1}},
  sen/.style     = {font=\scriptsize\ttfamily, inner sep=1.5pt},
  nodo/.style    = {font=\tiny\ttfamily, inner sep=1pt, black!75},
  pin/.style     = {font=\scriptsize, align=center, inner sep=1pt},
  grupo/.style   = {draw=black!55, dashed, line width=0.7pt, rounded corners=4pt},
  canal/.style   = {grupo},
  contenedor/.style = {draw=black!80, line width=1.0pt, rounded corners=6pt},
  % marca numerada: cada número es una fila de la tabla del documento
  marca/.style   = {circle, fill=black, text=white, font=\scriptsize\bfseries,
                    inner sep=0pt, minimum size=4.3mm},
  %
  % DECISIONES (docs/mca/vista_general_mca.md). Cada opción es un bloque, todas
  % entran a un mux (\muxdec, abajo) y el registro que elige entra por abajo.
  % la opción que trae el reset (el comportamiento histórico): borde doble
  defecto/.style = {double=white, double distance=0.8pt, line width=0.45pt},
  % rótulo de la decisión, el mismo Dn del panorama y del documento
  decision/.style = {circle, fill=black, text=white, font=\small\bfseries,
                     inner sep=0pt, minimum size=7mm},
  %
  % LÓGICA (docs/mca/figuras/contabilidad/ y los detalles de las decisiones)
  % registro: caja gris con el triángulo de reloj en el borde de abajo
  registro/.style = {draw=black, line width=0.6pt, fill=grisr, align=center,
                     font=\scriptsize, inner sep=2pt, minimum width=#1,
                     minimum height=0.85cm,
                     path picture={\draw[line width=0.5pt]
                       ([xshift=-1.1mm]path picture bounding box.south) --
                       ([yshift=1.6mm]path picture bounding box.south) --
                       ([xshift=1.1mm]path picture bounding box.south);}},
  registro/.default = 0.8cm,
  % condición de una cadena if / else-if: "sí" sale a la derecha, "no" abajo
  pregunta/.style = {draw=black, line width=0.5pt, rounded corners=7pt,
                     fill=white, font=\scriptsize\ttfamily, align=center,
                     inner sep=2.5pt, minimum height=0.62cm, text width=#1},
  pregunta/.default = 1.9cm,
  si/.style      = {font=\tiny, inner sep=1pt, above, pos=0.12},
  no/.style      = {font=\tiny, inner sep=1pt, left, pos=0.45},
  %
  titulo/.style  = {anchor=north west, font=\small, inner sep=4pt, align=left},
  rotulo/.style  = {anchor=north west, font=\scriptsize\itshape, black!70,
                    inner sep=2pt, align=left},
}

% submódulo: nombre en monoespaciada; lógica propia: nombre en romana
\newcommand{\blk}[2]{{\ttfamily\bfseries #1}\\[1pt]#2}
\newcommand{\etq}[2]{{\bfseries #1}\\[1pt]#2}
\newcommand{\punto}[1]{\fill (#1) circle (1.3pt);}

% Parámetros en el formato del datasheet: NOMBRE=valor en monoespaciada, una
% entrada por celda, hasta cuatro columnas. Una columna: \prm{A=1\\B=2};
% dos: \prm{A=1 & B=2\\C=3}.
\newcommand{\prm}[1]{{\ttfamily\fontsize{5.6}{6.6}\selectfont
  \begin{tabular}{@{}l@{\hspace{7pt}}l@{\hspace{7pt}}l@{\hspace{7pt}}l@{}}#1\end{tabular}}}

\def\anchopos{0.5}

% Mux de una decisión: trapecio con el lado ancho (las entradas) a la izquierda.
% #1 = nombre, #2 = centro, #3 = alto en cm. Deja tres coordenadas:
%   (#1-in)  centro del lado de entradas: las flechas llegan a (#1-in |- opcion.east)
%   (#1-out) la salida
%   (#1-sel) mitad del lado de abajo: por ahí entra el registro que elige
% El rótulo de adentro es opcional: \muxdec{m}{(4,0)}{3.6} y después
% \node[...] at (m) {...}.
\newcommand{\muxdec}[3]{%
  \path #2 coordinate (#1);
  \draw[line width=0.6pt, fill=grisr]
    ($(#1)+(-0.45,{#3/2})$) -- ($(#1)+(0.45,{#3/2-0.35})$) --
    ($(#1)+(0.45,{-#3/2+0.35})$) -- ($(#1)+(-0.45,{-#3/2})$) -- cycle;
  \coordinate (#1-in)  at ($(#1)+(-0.45,0)$);
  \coordinate (#1-out) at ($(#1)+(0.45,0)$);
  \coordinate (#1-sel) at ($(#1)+(0,{-#3/2+0.175})$);}

% El mismo mux rotado 90 grados, para un flujo de arriba hacia abajo: entradas
% por el lado ancho de arriba, salida abajo y el registro que elige por la
% izquierda. #1 = nombre, #2 = centro, #3 = ancho en cm. Deja (#1-in) centro del
% lado de entradas, (#1-out) la salida y (#1-sel) mitad del lado izquierdo.
\newcommand{\muxdecv}[3]{%
  \path #2 coordinate (#1);
  \draw[line width=0.6pt, fill=grisr]
    ($(#1)+({-#3/2},0.45)$) -- ($(#1)+({#3/2},0.45)$) --
    ($(#1)+({#3/2-0.35},-0.45)$) -- ($(#1)+({-#3/2+0.35},-0.45)$) -- cycle;
  \coordinate (#1-in)  at ($(#1)+(0,0.45)$);
  \coordinate (#1-out) at ($(#1)+(0,-0.45)$);
  \coordinate (#1-sel) at ($(#1)+({-#3/2+0.175},0)$);}

% Leyenda de colores + nota. #1 = posición, #2 = color/texto, #3 = nota.
\newcommand{\leyenda}[3]{%
  \begin{scope}[shift={(#1)}, every node/.style={font=\scriptsize}]
    \coordinate (t0) at (-0.5,0);
    \foreach \c/\t [count=\i] in {#2} {
      \pgfmathtruncatemacro\j{\i-1}
      \node[draw, fill=\c, minimum width=4mm, minimum height=3mm, inner sep=0pt,
            right=5mm of t\j.east, anchor=west] (l\i) {};
      \node[right=1.5mm of l\i] (t\i) {\t};
    }
    \node[anchor=north west, black!70, text width=19cm, align=left]
      at (-0.05,-0.3) {#3};
  \end{scope}}

% Lo que dicen los parámetros, para la nota de la leyenda
\newcommand{\notaparam}{\texttt{NOMBRE=valor}: parámetros con que el padre
  instancia el bloque (se omiten los que sólo pasan el del padre sin cambio).}
 (las figuras
% viven en docs/<módulo>/figuras/). Un solo lugar para la paleta, los estilos de
% bloque y el formato de los parámetros: si cambia algo acá, cambia en todas.
%
% La paleta es la misma de docs/figuras/estilo_svg.py y de PALETA en
% modulos/mk/datasheet.py; si se toca un color, se toca en los tres.
%
% Requiere: \usepackage{tikz} (lo da la clase standalone con la opción tikz).

\usetikzlibrary{arrows.meta,positioning,calc,fit,backgrounds,shapes.multipart,
                decorations.markings,shadows}

% ---------------------------------------------------------------- colores ---
\definecolor{azul}{RGB}{108,172,228}      % acondicionamiento
\definecolor{verde}{RGB}{110,214,142}     % captura / histogramas en BRAM
\definecolor{amarillo}{RGB}{252,232,160}  % trigger
\definecolor{naranja}{RGB}{252,206,160}   % camino a DDR
\definecolor{grisr}{RGB}{226,226,226}     % compartido / infraestructura
\definecolor{rojo}{RGB}{240,128,128}      % bus del sistema

% ---------------------------------------------------------------- estilos ---
\tikzset{
  >={Stealth[length=2.4mm,width=1.8mm]},
  % fondo blanco siempre: el SVG se ve igual en modo claro y oscuro
  fondo/.style   = {show background rectangle, inner frame sep=8pt,
                    background rectangle/.style={fill=white, draw=none}},
  bloque/.style  = {draw=black, line width=0.6pt, rounded corners=2pt,
                    align=center, text width=2.2cm, minimum height=1.15cm,
                    inner sep=2pt, font=\scriptsize, fill=#1},
  datos/.style   = {bloque=azul!35},
  trig/.style    = {bloque=amarillo},
  captura/.style = {bloque=verde!55},
  ddr/.style     = {bloque=naranja!70},
  comp/.style    = {bloque=grisr},
  bus/.style     = {bloque=rojo!45},
  % lógica escrita dentro del módulo, sin submódulo propio
  etapa/.style   = {bloque=#1, line width=0.5pt},
  etapa/.default = verde!20,
  % lo que el bitstream no trae (atado a cero, stub, salida muda)
  mudo/.style    = {bloque=white, draw=black!45, dashed, text=black!55},
  %
  % BLOQUE CON PARÁMETROS, en el formato del datasheet: el encabezado lleva el
  % nombre y los parámetros con que lo instancia el padre (NOMBRE=valor, en
  % monoespaciada, uno por línea); el cuerpo, más claro, lo que hace.
  %   \node[datos, param={azul!35}{azul!12}] {\blk{nombre}{} \prm{A=1\\B=2}
  %                                            \nodepart{two} función};
  param/.style 2 args = {rectangle split, rectangle split parts=2,
                         rectangle split part fill={#1, #2}},
  %
  flujo/.style   = {->, line width=0.8pt},
  ctrl/.style    = {->, line width=0.6pt, black!65},
  % la línea que pasa por encima en un cruce: un halo blanco la separa
  cruce/.style   = {preaction={draw=white, line width=3.4pt, -}},
  % marca de ancho de bus: barra oblicua a mitad de la flecha, con el número
  % encima. \draw[flujo, ancho=14] ...; para otra posición, ancho en=0.3.
  ancho/.style   = {decoration={markings, mark=at position \anchopos with {
                      \draw[line width=0.6pt, -] (-0.1,-0.14) -- (0.1,0.14);
                      \node[font=\scriptsize, above=1.5pt] {#1};}},
                    postaction={decorate}},
  % la misma marca en una flecha vertical: el número va a la derecha
  anchov/.style  = {decoration={markings, mark=at position \anchopos with {
                      \draw[line width=0.6pt, -] (-0.14,-0.1) -- (0.14,0.1);
                      \node[font=\scriptsize, right=3pt] {#1};}},
                    postaction={decorate}},
  ancho en/.code = {\def\anchopos{#1}},
  sen/.style     = {font=\scriptsize\ttfamily, inner sep=1.5pt},
  nodo/.style    = {font=\tiny\ttfamily, inner sep=1pt, black!75},
  pin/.style     = {font=\scriptsize, align=center, inner sep=1pt},
  grupo/.style   = {draw=black!55, dashed, line width=0.7pt, rounded corners=4pt},
  canal/.style   = {grupo},
  contenedor/.style = {draw=black!80, line width=1.0pt, rounded corners=6pt},
  % marca numerada: cada número es una fila de la tabla del documento
  marca/.style   = {circle, fill=black, text=white, font=\scriptsize\bfseries,
                    inner sep=0pt, minimum size=4.3mm},
  %
  % DECISIONES (docs/mca/vista_general_mca.md). Cada opción es un bloque, todas
  % entran a un mux (\muxdec, abajo) y el registro que elige entra por abajo.
  % la opción que trae el reset (el comportamiento histórico): borde doble
  defecto/.style = {double=white, double distance=0.8pt, line width=0.45pt},
  % rótulo de la decisión, el mismo Dn del panorama y del documento
  decision/.style = {circle, fill=black, text=white, font=\small\bfseries,
                     inner sep=0pt, minimum size=7mm},
  %
  % LÓGICA (docs/mca/figuras/contabilidad/ y los detalles de las decisiones)
  % registro: caja gris con el triángulo de reloj en el borde de abajo
  registro/.style = {draw=black, line width=0.6pt, fill=grisr, align=center,
                     font=\scriptsize, inner sep=2pt, minimum width=#1,
                     minimum height=0.85cm,
                     path picture={\draw[line width=0.5pt]
                       ([xshift=-1.1mm]path picture bounding box.south) --
                       ([yshift=1.6mm]path picture bounding box.south) --
                       ([xshift=1.1mm]path picture bounding box.south);}},
  registro/.default = 0.8cm,
  % condición de una cadena if / else-if: "sí" sale a la derecha, "no" abajo
  pregunta/.style = {draw=black, line width=0.5pt, rounded corners=7pt,
                     fill=white, font=\scriptsize\ttfamily, align=center,
                     inner sep=2.5pt, minimum height=0.62cm, text width=#1},
  pregunta/.default = 1.9cm,
  si/.style      = {font=\tiny, inner sep=1pt, above, pos=0.12},
  no/.style      = {font=\tiny, inner sep=1pt, left, pos=0.45},
  %
  titulo/.style  = {anchor=north west, font=\small, inner sep=4pt, align=left},
  rotulo/.style  = {anchor=north west, font=\scriptsize\itshape, black!70,
                    inner sep=2pt, align=left},
}

% submódulo: nombre en monoespaciada; lógica propia: nombre en romana
\newcommand{\blk}[2]{{\ttfamily\bfseries #1}\\[1pt]#2}
\newcommand{\etq}[2]{{\bfseries #1}\\[1pt]#2}
\newcommand{\punto}[1]{\fill (#1) circle (1.3pt);}

% Parámetros en el formato del datasheet: NOMBRE=valor en monoespaciada, una
% entrada por celda, hasta cuatro columnas. Una columna: \prm{A=1\\B=2};
% dos: \prm{A=1 & B=2\\C=3}.
\newcommand{\prm}[1]{{\ttfamily\fontsize{5.6}{6.6}\selectfont
  \begin{tabular}{@{}l@{\hspace{7pt}}l@{\hspace{7pt}}l@{\hspace{7pt}}l@{}}#1\end{tabular}}}

\def\anchopos{0.5}

% Mux de una decisión: trapecio con el lado ancho (las entradas) a la izquierda.
% #1 = nombre, #2 = centro, #3 = alto en cm. Deja tres coordenadas:
%   (#1-in)  centro del lado de entradas: las flechas llegan a (#1-in |- opcion.east)
%   (#1-out) la salida
%   (#1-sel) mitad del lado de abajo: por ahí entra el registro que elige
% El rótulo de adentro es opcional: \muxdec{m}{(4,0)}{3.6} y después
% \node[...] at (m) {...}.
\newcommand{\muxdec}[3]{%
  \path #2 coordinate (#1);
  \draw[line width=0.6pt, fill=grisr]
    ($(#1)+(-0.45,{#3/2})$) -- ($(#1)+(0.45,{#3/2-0.35})$) --
    ($(#1)+(0.45,{-#3/2+0.35})$) -- ($(#1)+(-0.45,{-#3/2})$) -- cycle;
  \coordinate (#1-in)  at ($(#1)+(-0.45,0)$);
  \coordinate (#1-out) at ($(#1)+(0.45,0)$);
  \coordinate (#1-sel) at ($(#1)+(0,{-#3/2+0.175})$);}

% El mismo mux rotado 90 grados, para un flujo de arriba hacia abajo: entradas
% por el lado ancho de arriba, salida abajo y el registro que elige por la
% izquierda. #1 = nombre, #2 = centro, #3 = ancho en cm. Deja (#1-in) centro del
% lado de entradas, (#1-out) la salida y (#1-sel) mitad del lado izquierdo.
\newcommand{\muxdecv}[3]{%
  \path #2 coordinate (#1);
  \draw[line width=0.6pt, fill=grisr]
    ($(#1)+({-#3/2},0.45)$) -- ($(#1)+({#3/2},0.45)$) --
    ($(#1)+({#3/2-0.35},-0.45)$) -- ($(#1)+({-#3/2+0.35},-0.45)$) -- cycle;
  \coordinate (#1-in)  at ($(#1)+(0,0.45)$);
  \coordinate (#1-out) at ($(#1)+(0,-0.45)$);
  \coordinate (#1-sel) at ($(#1)+({-#3/2+0.175},0)$);}

% Leyenda de colores + nota. #1 = posición, #2 = color/texto, #3 = nota.
\newcommand{\leyenda}[3]{%
  \begin{scope}[shift={(#1)}, every node/.style={font=\scriptsize}]
    \coordinate (t0) at (-0.5,0);
    \foreach \c/\t [count=\i] in {#2} {
      \pgfmathtruncatemacro\j{\i-1}
      \node[draw, fill=\c, minimum width=4mm, minimum height=3mm, inner sep=0pt,
            right=5mm of t\j.east, anchor=west] (l\i) {};
      \node[right=1.5mm of l\i] (t\i) {\t};
    }
    \node[anchor=north west, black!70, text width=19cm, align=left]
      at (-0.05,-0.3) {#3};
  \end{scope}}

% Lo que dicen los parámetros, para la nota de la leyenda
\newcommand{\notaparam}{\texttt{NOMBRE=valor}: parámetros con que el padre
  instancia el bloque (se omiten los que sólo pasan el del padre sin cambio).}
 (las figuras
% viven en docs/<módulo>/figuras/). Un solo lugar para la paleta, los estilos de
% bloque y el formato de los parámetros: si cambia algo acá, cambia en todas.
%
% La paleta es la misma de docs/figuras/estilo_svg.py y de PALETA en
% modulos/mk/datasheet.py; si se toca un color, se toca en los tres.
%
% Requiere: \usepackage{tikz} (lo da la clase standalone con la opción tikz).

\usetikzlibrary{arrows.meta,positioning,calc,fit,backgrounds,shapes.multipart,
                decorations.markings,shadows}

% ---------------------------------------------------------------- colores ---
\definecolor{azul}{RGB}{108,172,228}      % acondicionamiento
\definecolor{verde}{RGB}{110,214,142}     % captura / histogramas en BRAM
\definecolor{amarillo}{RGB}{252,232,160}  % trigger
\definecolor{naranja}{RGB}{252,206,160}   % camino a DDR
\definecolor{grisr}{RGB}{226,226,226}     % compartido / infraestructura
\definecolor{rojo}{RGB}{240,128,128}      % bus del sistema

% ---------------------------------------------------------------- estilos ---
\tikzset{
  >={Stealth[length=2.4mm,width=1.8mm]},
  % fondo blanco siempre: el SVG se ve igual en modo claro y oscuro
  fondo/.style   = {show background rectangle, inner frame sep=8pt,
                    background rectangle/.style={fill=white, draw=none}},
  bloque/.style  = {draw=black, line width=0.6pt, rounded corners=2pt,
                    align=center, text width=2.2cm, minimum height=1.15cm,
                    inner sep=2pt, font=\scriptsize, fill=#1},
  datos/.style   = {bloque=azul!35},
  trig/.style    = {bloque=amarillo},
  captura/.style = {bloque=verde!55},
  ddr/.style     = {bloque=naranja!70},
  comp/.style    = {bloque=grisr},
  bus/.style     = {bloque=rojo!45},
  % lógica escrita dentro del módulo, sin submódulo propio
  etapa/.style   = {bloque=#1, line width=0.5pt},
  etapa/.default = verde!20,
  % lo que el bitstream no trae (atado a cero, stub, salida muda)
  mudo/.style    = {bloque=white, draw=black!45, dashed, text=black!55},
  %
  % BLOQUE CON PARÁMETROS, en el formato del datasheet: el encabezado lleva el
  % nombre y los parámetros con que lo instancia el padre (NOMBRE=valor, en
  % monoespaciada, uno por línea); el cuerpo, más claro, lo que hace.
  %   \node[datos, param={azul!35}{azul!12}] {\blk{nombre}{} \prm{A=1\\B=2}
  %                                            \nodepart{two} función};
  param/.style 2 args = {rectangle split, rectangle split parts=2,
                         rectangle split part fill={#1, #2}},
  %
  flujo/.style   = {->, line width=0.8pt},
  ctrl/.style    = {->, line width=0.6pt, black!65},
  % la línea que pasa por encima en un cruce: un halo blanco la separa
  cruce/.style   = {preaction={draw=white, line width=3.4pt, -}},
  % marca de ancho de bus: barra oblicua a mitad de la flecha, con el número
  % encima. \draw[flujo, ancho=14] ...; para otra posición, ancho en=0.3.
  ancho/.style   = {decoration={markings, mark=at position \anchopos with {
                      \draw[line width=0.6pt, -] (-0.1,-0.14) -- (0.1,0.14);
                      \node[font=\scriptsize, above=1.5pt] {#1};}},
                    postaction={decorate}},
  % la misma marca en una flecha vertical: el número va a la derecha
  anchov/.style  = {decoration={markings, mark=at position \anchopos with {
                      \draw[line width=0.6pt, -] (-0.14,-0.1) -- (0.14,0.1);
                      \node[font=\scriptsize, right=3pt] {#1};}},
                    postaction={decorate}},
  ancho en/.code = {\def\anchopos{#1}},
  sen/.style     = {font=\scriptsize\ttfamily, inner sep=1.5pt},
  nodo/.style    = {font=\tiny\ttfamily, inner sep=1pt, black!75},
  pin/.style     = {font=\scriptsize, align=center, inner sep=1pt},
  grupo/.style   = {draw=black!55, dashed, line width=0.7pt, rounded corners=4pt},
  canal/.style   = {grupo},
  contenedor/.style = {draw=black!80, line width=1.0pt, rounded corners=6pt},
  % marca numerada: cada número es una fila de la tabla del documento
  marca/.style   = {circle, fill=black, text=white, font=\scriptsize\bfseries,
                    inner sep=0pt, minimum size=4.3mm},
  %
  % DECISIONES (docs/mca/vista_general_mca.md). Cada opción es un bloque, todas
  % entran a un mux (\muxdec, abajo) y el registro que elige entra por abajo.
  % la opción que trae el reset (el comportamiento histórico): borde doble
  defecto/.style = {double=white, double distance=0.8pt, line width=0.45pt},
  % rótulo de la decisión, el mismo Dn del panorama y del documento
  decision/.style = {circle, fill=black, text=white, font=\small\bfseries,
                     inner sep=0pt, minimum size=7mm},
  %
  % LÓGICA (docs/mca/figuras/contabilidad/ y los detalles de las decisiones)
  % registro: caja gris con el triángulo de reloj en el borde de abajo
  registro/.style = {draw=black, line width=0.6pt, fill=grisr, align=center,
                     font=\scriptsize, inner sep=2pt, minimum width=#1,
                     minimum height=0.85cm,
                     path picture={\draw[line width=0.5pt]
                       ([xshift=-1.1mm]path picture bounding box.south) --
                       ([yshift=1.6mm]path picture bounding box.south) --
                       ([xshift=1.1mm]path picture bounding box.south);}},
  registro/.default = 0.8cm,
  % condición de una cadena if / else-if: "sí" sale a la derecha, "no" abajo
  pregunta/.style = {draw=black, line width=0.5pt, rounded corners=7pt,
                     fill=white, font=\scriptsize\ttfamily, align=center,
                     inner sep=2.5pt, minimum height=0.62cm, text width=#1},
  pregunta/.default = 1.9cm,
  si/.style      = {font=\tiny, inner sep=1pt, above, pos=0.12},
  no/.style      = {font=\tiny, inner sep=1pt, left, pos=0.45},
  %
  titulo/.style  = {anchor=north west, font=\small, inner sep=4pt, align=left},
  rotulo/.style  = {anchor=north west, font=\scriptsize\itshape, black!70,
                    inner sep=2pt, align=left},
}

% submódulo: nombre en monoespaciada; lógica propia: nombre en romana
\newcommand{\blk}[2]{{\ttfamily\bfseries #1}\\[1pt]#2}
\newcommand{\etq}[2]{{\bfseries #1}\\[1pt]#2}
\newcommand{\punto}[1]{\fill (#1) circle (1.3pt);}

% Parámetros en el formato del datasheet: NOMBRE=valor en monoespaciada, una
% entrada por celda, hasta cuatro columnas. Una columna: \prm{A=1\\B=2};
% dos: \prm{A=1 & B=2\\C=3}.
\newcommand{\prm}[1]{{\ttfamily\fontsize{5.6}{6.6}\selectfont
  \begin{tabular}{@{}l@{\hspace{7pt}}l@{\hspace{7pt}}l@{\hspace{7pt}}l@{}}#1\end{tabular}}}

\def\anchopos{0.5}

% Mux de una decisión: trapecio con el lado ancho (las entradas) a la izquierda.
% #1 = nombre, #2 = centro, #3 = alto en cm. Deja tres coordenadas:
%   (#1-in)  centro del lado de entradas: las flechas llegan a (#1-in |- opcion.east)
%   (#1-out) la salida
%   (#1-sel) mitad del lado de abajo: por ahí entra el registro que elige
% El rótulo de adentro es opcional: \muxdec{m}{(4,0)}{3.6} y después
% \node[...] at (m) {...}.
\newcommand{\muxdec}[3]{%
  \path #2 coordinate (#1);
  \draw[line width=0.6pt, fill=grisr]
    ($(#1)+(-0.45,{#3/2})$) -- ($(#1)+(0.45,{#3/2-0.35})$) --
    ($(#1)+(0.45,{-#3/2+0.35})$) -- ($(#1)+(-0.45,{-#3/2})$) -- cycle;
  \coordinate (#1-in)  at ($(#1)+(-0.45,0)$);
  \coordinate (#1-out) at ($(#1)+(0.45,0)$);
  \coordinate (#1-sel) at ($(#1)+(0,{-#3/2+0.175})$);}

% El mismo mux rotado 90 grados, para un flujo de arriba hacia abajo: entradas
% por el lado ancho de arriba, salida abajo y el registro que elige por la
% izquierda. #1 = nombre, #2 = centro, #3 = ancho en cm. Deja (#1-in) centro del
% lado de entradas, (#1-out) la salida y (#1-sel) mitad del lado izquierdo.
\newcommand{\muxdecv}[3]{%
  \path #2 coordinate (#1);
  \draw[line width=0.6pt, fill=grisr]
    ($(#1)+({-#3/2},0.45)$) -- ($(#1)+({#3/2},0.45)$) --
    ($(#1)+({#3/2-0.35},-0.45)$) -- ($(#1)+({-#3/2+0.35},-0.45)$) -- cycle;
  \coordinate (#1-in)  at ($(#1)+(0,0.45)$);
  \coordinate (#1-out) at ($(#1)+(0,-0.45)$);
  \coordinate (#1-sel) at ($(#1)+({-#3/2+0.175},0)$);}

% Leyenda de colores + nota. #1 = posición, #2 = color/texto, #3 = nota.
\newcommand{\leyenda}[3]{%
  \begin{scope}[shift={(#1)}, every node/.style={font=\scriptsize}]
    \coordinate (t0) at (-0.5,0);
    \foreach \c/\t [count=\i] in {#2} {
      \pgfmathtruncatemacro\j{\i-1}
      \node[draw, fill=\c, minimum width=4mm, minimum height=3mm, inner sep=0pt,
            right=5mm of t\j.east, anchor=west] (l\i) {};
      \node[right=1.5mm of l\i] (t\i) {\t};
    }
    \node[anchor=north west, black!70, text width=19cm, align=left]
      at (-0.05,-0.3) {#3};
  \end{scope}}

% Lo que dicen los parámetros, para la nota de la leyenda
\newcommand{\notaparam}{\texttt{NOMBRE=valor}: parámetros con que el padre
  instancia el bloque (se omiten los que sólo pasan el del padre sin cambio).}

\tikzset{
  op/.style      = {oper, font=\scriptsize},
  caja/.style    = {desp, inner sep=2pt, minimum height=0.5cm},
  reg/.style     = {registro=#1, minimum height=0.55cm, font=\scriptsize\ttfamily},
  reg/.default   = 1.15cm,
  fondo f/.style = {draw=black!45, dashed, line width=0.5pt, rounded corners=3pt, fill=#1},
  rotf/.style    = {font=\scriptsize\itshape, black!75, inner sep=2pt},
  red/.style     = {font=\scriptsize\ttfamily, inner sep=1pt},
  chico/.style   = {font=\tiny, inner sep=1pt, align=center},
  y3/.style      = {logic gate input sep=0.2cm},
  y2/.style      = {logic gate input sep=0.2cm},
}
% \red[largo]{pin}{nombre}: la red que entra a un pin, por nombre, desde la izquierda
\newcommand{\red}[3][0.3]{\draw[flujo] ($(#2)+(-#1,0)$) node[left, red] {#3} -- (#2);}
% \redab[largo]{pin}{nombre}: la misma, desde abajo
\newcommand{\redab}[3][0.3]{\draw[flujo] ($(#2)+(0,-#1)$) node[below, red] {#3} -- (#2);}
% \redar[largo]{pin}{nombre}: la misma, desde arriba
\newcommand{\redar}[3][0.3]{\draw[flujo] ($(#2)+(0,#1)$) node[above, red] {#3} -- (#2);}
% \redsal[largo]{pin}{nombre}: la red que sale de un pin hacia la derecha
\newcommand{\redsal}[3][0.3]{\draw[flujo] (#2) -- ($(#2)+(#1,0)$) node[right, red] {#3};}
% \sal{pin}{puerto}: la salida, hasta el borde
\newcommand{\sal}[2]{\draw[flujo] (#1) -- (\xout,0 |- #1) node[right, sen] {#2};}
% \mux2{nombre}{centro}{alto}: el mux de carga, con 0 arriba y 1 abajo
\newcommand{\muxdos}[3]{\muxdec{#1}{#2}{#3}
  \node[chico, anchor=west] at ($(#1-in)+(0.04,{#3/2-0.25})$) {0};
  \node[chico, anchor=west] at ($(#1-in)+(0.04,{-#3/2+0.25})$) {1};}

\begin{document}
\begin{tikzpicture}[fondo]

\def\xin{-0.3}    % donde arrancan las entradas
\def\xout{19.75}  % donde terminan las salidas
\def\ytop{15.45}  % de donde bajan las de configuración

% ============================================================ las regiones ===
\node[fondo f=amarillo!22, fit={(0.35,10.8) (11.9,14.65)}] (fcmp) {};
\node[fondo f=amarillo!22, fit={(12.3,10.8) (19.05,14.65)}] (fest) {};
\node[fondo f=verde!14, fit={(0.35,3.95) (11.9,10.4)}] (fven) {};
\node[fondo f=verde!14, fit={(12.3,8.3) (19.05,10.4)}] (fpic) {};
\node[fondo f=verde!14, fit={(12.3,5.55) (19.05,7.9)}] (fcar) {};
\node[fondo f=verde!14, fit={(12.3,3.95) (19.05,5.3)}] (famp) {};
\node[fondo f=verde!14, fit={(0.35,1.6) (19.05,3.55)}] (fmes) {};
\node[fondo f=grisr!45, fit={(0.35,-3.3) (19.05,1.2)}] (fmar) {};
\node[rotf, anchor=north west] at (fcmp.north west) {compuertas};
\node[rotf, anchor=north west] at (fest.north west) {estado};
\node[rotf, anchor=north east] at (fven.north east) {ventana de amplitud};
\node[rotf, anchor=north west] at (fpic.north west) {pico};
\node[rotf, anchor=north west] at (fcar.north west) {cargas};
\node[rotf, anchor=north west] at (famp.north west) {amplitud a publicar};
\node[rotf, anchor=north east] at (fmes.north east) {meseta, una por forma ($\times$3)};
\node[rotf, anchor=north west] at (fmar.north west) {marcas de solape};

% ============================================================ compuertas ===
% abre: el asignador eligió esta ranura y la muestra es válida
\node[puerta y, y2, anchor=input 1] (gab) at (1.15,13.95) {};
\draw[flujo] (\xin,0 |- gab.input 1) node[left, sen] {abre\_i} -- (gab.input 1);
\draw[flujo] (\xin,0 |- gab.input 2) node[left, sen] {val\_i} -- (gab.input 2);
\redsal{gab.output}{abre}
% len: muestras desde el disparo
\node[reg=0.9cm] (rlen) at (3.6,13.75) {len};
\node[op] (lmas) at (4.85,13.75) {$+1$};
\draw[flujo] (rlen.east) -- (lmas.west);
\draw[flujo] (lmas.north) |- ($(rlen.north)+(0,0.3)$) -- (rlen.north);
\node[chico, anchor=south west] at ($(rlen.north)+(0.1,0.32)$) {1 con \texttt{abre}; $+1$ con \texttt{acumula}};
% las compuertas, retenidas al abrir
\node[reg=1.0cm] (rgs) at (3.6,12.65) {g\_short};
\node[reg=1.0cm] (rgl) at (3.6,11.65) {g\_long};
\draw[ctrl] (\xin,0 |- rgs.west) node[left, sen, black!80] {cfg\_gate\_short\_i} -- (rgs.west);
\draw[ctrl] (\xin,0 |- rgl.west) node[left, sen, black!80] {cfg\_gate\_long\_i} -- (rgl.west);
\node[chico, anchor=north] at (rgl.south) {se cargan con \texttt{abre}};
\node[op] (ges) at (4.85,12.65) {$\geq$};
\node[op] (gel) at (4.85,11.65) {$\geq$};
\draw[flujo] (rgs.east) -- (ges.west);
\draw[flujo] (rgl.east) -- (gel.west);
\redar{ges.north}{len}
\redab{gel.south}{len}
% cierre, acumula, in_tail y estancada
\node[puerta y=nnn, y3, anchor=input 1] (gci) at (7.7,14.3) {};
\node[puerta y=nni, y3, anchor=input 1] (gac) at (7.7,13.25) {};
\node[puerta y, y2, anchor=input 1] (git) at (7.7,12.3) {};
\node[puerta y=nin, y3, anchor=input 1] (gst) at (7.7,11.6) {};
\red[0.45]{gci.input 1}{abierta}
\red[0.45]{gci.input 2}{val\_i}
\red[0.45]{gac.input 1}{abierta}
\red[0.45]{gac.input 2}{val\_i}
\red[0.45]{gac.input 3}{cierre}
\red[0.45]{git.input 1}{abierta}
\red[0.45]{gst.input 1}{abierta}
\red[0.45]{gst.input 2}{val\_i}
\draw[flujo] (gel.east) -- (5.75,0 |- gel.east) coordinate (jgl);
\draw[flujo] (jgl) |- (gci.input 3);
\draw[flujo] (jgl) |- (gst.input 3);
\draw[flujo, cruce] (ges.east) -- (6.1,0 |- ges.east) |- (git.input 2);
\redsal{gci.output}{cierre}
\redsal{gac.output}{acumula}
\redsal{git.output}{in\_tail}
\redsal{gst.output}{estancada}

\begin{scope}[xshift=0.9cm]\def\xout{18.85}
% ================================================================ estado ===
\node[registro=1.0cm, minimum height=1.5cm, font=\scriptsize\ttfamily] (rest) at (12.9,13.25) {est};
\red{rest.west |- 0,13.7}{abre}
\red{rest.west |- 0,13.25}{cierre}
\red{rest.west |- 0,12.8}{amp\_ok}
\draw[ctrl] (12.9,\ytop) node[above, sen, black!80] {libera\_i} -- (rest.north);
\node[chico, anchor=north west, align=left] at (11.55,12.3)
  {\textsc{libre} $\to$ \textsc{abierta}: \texttt{abre}\\
   \textsc{abierta} $\to$ \textsc{cerrada}: \texttt{cierre} $\wedge$ \texttt{amp\_ok}\\
   \textsc{abierta} $\to$ \textsc{libre}: \texttt{cierre} $\wedge\,\neg$\texttt{amp\_ok}\\
   \textsc{cerrada} $\to$ \textsc{libre}: \texttt{libera\_i}};
\node[caja, minimum width=1.5cm] (dli) at (14.9,13.8) {$=$ \textsc{libre}};
\node[caja, minimum width=1.5cm] (dab) at (14.9,13.25) {$=$ \textsc{abierta}};
\node[caja, minimum width=1.5cm] (dce) at (14.9,12.7) {$=$ \textsc{cerrada}};
\foreach \d in {dli, dab, dce} \draw[flujo] (rest.east |- \d.west) -- (\d.west);
\sal{dli.east}{libre\_o}
\sal{dab.east}{abierta\_o}
\node[chico, anchor=south] at ($(dab.east)+(0.45,0.02)$) {abierta};
\sal{dce.east}{cerrada\_o}
\node[puerta y, y2, anchor=input 1] (gacp) at (16.75,11.95) {};
\red{gacp.input 1}{cierre}
\red{gacp.input 2}{amp\_ok}
\sal{gacp.output}{acepta\_o}
\draw[flujo] (16.6,11.3) node[left, red] {cierre} -- (\xout,11.3) node[right, sen] {cierra\_o};
\draw[flujo] (16.6,10.95) node[left, red] {estancada} -- (\xout,10.95) node[right, sen] {estancada\_o};
\draw[flujo] (16.6,14.45) node[left, red] {len} -- (\xout,14.45) node[right, sen] {len\_o};

\end{scope}

% ==================================================== ventana de amplitud ===
% el pico que quedaría con esta muestra, contra [min, max]
\draw[flujo] (1.25,9.85) node[left, red] {xc\_i} -- (3.05,9.85);
\node[op] (pmx) at (2.2,9.25) {max};
\draw[flujo] (2.2,9.85) -- (pmx.north);
\punto{2.2,9.85}
\red[0.45]{pmx.west}{peak}
\muxdos{mc}{(3.5,9.55)}{1.05}
\draw[flujo] (pmx.east) -- (mc-in |- pmx.east);
\redab[0.2]{mc-sel}{abierta}
\draw[flujo] (mc-out) -- (4.2,0 |- mc-out) coordinate (jmc);
\node[op] (cmn) at (4.85,9.85) {$\geq$};
\node[op] (cmx) at (4.85,9.2) {$\leq$};
\draw[flujo] (jmc) |- (cmn.west);
\draw[flujo] (jmc) |- (cmx.west);
\punto{jmc}
\redar[0.2]{cmn.north}{cfg\_amp\_min\_i}
\redab[0.2]{cmx.south}{cfg\_amp\_max\_i}
\node[puerta y, y2, anchor=input 1] (gpk) at (5.95,9.665) {};
\draw[flujo] (cmn.east) -| ($(gpk.input 1)+(-0.2,0)$) -- (gpk.input 1);
\draw[flujo] (cmx.east) -| ($(gpk.input 2)+(-0.2,0)$) -- (gpk.input 2);
\node[reg=1.2cm] (rpk) at (7.75,0 |- gpk.output) {peak\_ok\_r};
\draw[flujo] (gpk.output) -- (rpk.west);

% sombra del mínimo: d_min = Q - min_q, con Q la carga hasta esta muestra
\node[op] (s1) at (2.2,8.0) {$+$};
\node[op] (s2) at (2.2,7.0) {$+$};
\red[0.45]{s1.west}{xc\_i}
\redar[0.2]{s1.north}{$-$min\_q\_i}
\red[0.45]{s2.west}{d\_min}
\draw[flujo] ($(s2.south)+(0,-0.2)$) node[right, red] {xc\_i} -- (s2.south);
\muxdos{ma}{(3.5,7.5)}{1.55}
\draw[flujo] (s1.east) -- (ma-in |- s1.east);
\draw[flujo] (s2.east) -- (ma-in |- s2.east);
\redab{ma-sel}{abierta}
\node[reg=0.95cm] (rdmin) at (5.2,8.0) {d\_min};
\draw[flujo] (ma-out) -- (4.2,0 |- ma-out) coordinate (jma) |- (rdmin.west);
\redsal[0.2]{rdmin.east}{d\_min}
\node[op] (sga) at (5.2,7.05) {$\geq0$};
\draw[flujo] (jma) |- (sga.west);
\punto{jma}
\node[puerta y=ni, y2, anchor=input 1] (gfu) at (5.85,0 |- sga.east) {};
\draw[flujo] (sga.east) -- (gfu.input 1);
\redab[0.3]{gfu.input 2}{min\_fuera\_i}
\node[reg=1.2cm] (rgm) at (7.75,0 |- gfu.input 1) {ge\_min\_r};
\draw[flujo] (gfu.output) -- (rgm.west |- gfu.output);

% sombra del máximo: d_max = max_q - Q
\node[op] (s3) at (2.2,5.6) {$+$};
\node[op] (s4) at (2.2,4.6) {$+$};
\red[0.45]{s3.west}{max\_q\_i}
\draw[flujo] ($(s3.north)+(0,0.2)$) node[right, red] {$-$xc\_i} -- (s3.north);
\red[0.45]{s4.west}{d\_max}
\redab[0.2]{s4.south}{$-$xc\_i}
\muxdos{mb}{(3.5,5.1)}{1.55}
\draw[flujo] (s3.east) -- (mb-in |- s3.east);
\draw[flujo] (s4.east) -- (mb-in |- s4.east);
\redab{mb-sel}{abierta}
\node[reg=0.95cm] (rdmax) at (5.2,5.6) {d\_max};
\draw[flujo] (mb-out) -- (4.2,0 |- mb-out) coordinate (jmb) |- (rdmax.west);
\redsal[0.2]{rdmax.east}{d\_max}
\node[op] (sgb) at (5.2,4.65) {$\geq0$};
\draw[flujo] (jmb) |- (sgb.west);
\punto{jmb}
\node[reg=1.2cm] (rlm) at (7.75,4.65) {le\_max\_r};
\draw[flujo] (sgb.east) -- (rlm.west);
\node[chico, anchor=north west, align=left] at (6.3,6.25)
  {\texttt{d\_min}, \texttt{d\_max}: en \texttt{abre} $\vee$ \texttt{acumula}\\
   $x_c$ suma sólo con \texttt{val\_i}\\ los \texttt{*\_r}, en cada ciclo};

% amp_ok: por el pico, su ventana; por la integral, las dos sombras (o el techo saturado)
\node[puerta o, y2, anchor=input 2] (gsa) at (8.6,0 |- rlm.east) {};
\draw[flujo] (rlm.east) -- (gsa.input 2);
\draw[flujo] ($(gsa.input 1)+(-0.15,0.35)$) node[above, red] {sat\_ok\_i} |- (gsa.input 1);
\node[puerta y, y2, anchor=input 1] (gin) at (9.45,0 |- rgm.east) {};
\draw[flujo] (rgm.east) -- (gin.input 1);
\draw[flujo] (gsa.output) -| ($(gin.input 2)+(-0.15,0)$) -- (gin.input 2);
\muxdos{ms}{(10.95,7.75)}{1.6}
\draw[flujo] (rpk.east) -| ($(ms-in)+(-0.25,0.55)$) -- (ms-in |- 0,8.3);
\draw[flujo] (gin.output) -- (ms-in |- gin.output);
\redab[0.85]{ms-sel}{cfg\_amp\_src\_i}
\draw[flujo] (ms-out) -- ++(0.2,0) |- ++(0.15,0.75) node[above left, red, xshift=2pt] {amp\_ok};

\begin{scope}[xshift=0.9cm]\def\xout{18.85}
% =================================================================== pico ===
\node[op] (gpt) at (12.5,9.45) {$>$};
\red{gpt.west}{xc\_i}
\redab[0.25]{gpt.south}{peak}
\node[puerta y, y2, anchor=input 2] (gnp) at ($(gpt.east)+(0.3,0)$) {};
\draw[flujo] (gpt.east) -- (gnp.input 2);
\draw[flujo] ($(gnp.input 1)+(-0.25,0.35)$) node[above, red] {acumula} |- (gnp.input 1);
\redsal{gnp.output}{nuevo}
\node[reg=1.1cm] (rpi) at (16.0,9.8) {peak};
\node[reg=1.1cm] (rtp) at (16.0,9.0) {t\_peak};
\red{rpi.west}{xc\_i}
\red{rtp.west}{len}
\node[chico, anchor=south] at ($(rpi.north)+(0,0.03)$) {en: \texttt{abre} $\vee$ \texttt{nuevo}};
\node[chico, anchor=north] at ($(rtp.south)+(0,-0.03)$) {en: \texttt{nuevo}; 0 con \texttt{abre}};
\sal{rpi.east}{peak\_o}
\node[chico, anchor=south] at ($(rpi.east)+(0.5,0.02)$) {peak};
\sal{rtp.east}{t\_peak\_o}

% ================================================================= cargas ===
\node[op] (qa) at (12.9,7.35) {$+$};
\red{qa.west}{xc\_i}
\redar[0.2]{qa.north}{q\_tot}
\muxdos{mq}{(13.9,7.1)}{1.0}
\draw[flujo] (qa.east) -- (mq-in |- qa.east);
\draw[flujo] (13.25,6.85) node[left, red] {xc\_i} -- (mq-in |- 0,6.85);
\redab[0.2]{mq-sel}{abre}
\node[reg=1.1cm] (rqt) at (16.0,0 |- mq-out) {q\_tot};
\draw[flujo] (mq-out) -- (rqt.west);
\node[chico, anchor=south] at ($(rqt.north)+(0,0.03)$) {en: \texttt{abre} $\vee$ \texttt{acumula}};
\sal{rqt.east}{q\_tot\_o}
\node[chico, anchor=south] at ($(rqt.east)+(0.5,0.02)$) {q\_tot};
\node[op] (qb) at (12.9,6.35) {$+$};
\red{qb.west}{xc\_i}
\redab[0.2]{qb.south}{q\_tail}
\node[reg=1.1cm] (rqc) at (16.0,6.35) {q\_tail};
\draw[flujo] (qb.east) -- (rqc.west);
\node[chico, anchor=north] at ($(rqc.south)+(0,-0.03)$) {en: \texttt{abre} $\vee$ (\texttt{acumula} $\wedge$ \texttt{in\_tail}); 0 con \texttt{abre}};
\sal{rqc.east}{q\_tail\_o}

% ===================================================== amplitud a publicar ===
\node[caja] (qsh) at (12.6,4.75) {$\gg$ s};
\red{qsh.west}{q\_tot}
\redab[0.2]{qsh.south}{cfg\_q\_shift\_i}
\node[satur] (sat) at (13.7,4.75) {};
\draw[flujo] (qsh.east) -- (sat.west);
\node[chico, anchor=north] at (sat.south) {a \texttt{0xFFFF}};
\draw[flujo] (sat.east) -- (14.5,0 |- sat.east) coordinate (jsa);
\muxdos{mp}{(16.0,4.75)}{1.0}
\draw[flujo] (jsa) |- (mp-in |- 0,4.5);
\punto{jsa}
\draw[flujo] (15.35,5.0) node[left, red] {peak} -- (mp-in |- 0,5.0);
\draw[flujo] ($(mp-sel)+(0.3,-0.3)$) node[right, red] {cfg\_amp\_src\_i} -| (mp-sel);
\sal{mp-out}{amp\_sel\_o}
\draw[flujo] (jsa) |- (\xout,4.25) node[right, sen] {amp\_int\_o};

\end{scope}

% ======================================================= meseta, por forma ===
\node[op] (mi) at (2.6,2.95) {$\geq$};
\node[op] (mf) at (2.6,2.2) {$<$};
\red{mi.west}{len}
\red{mf.west}{len}
\redar[0.2]{mi.north}{}
\node[red, anchor=west] at ($(mi.north)+(0.08,0.12)$) {trap\_ini\_i};
\redab[0.2]{mf.south}{}
\node[red, anchor=west] at ($(mf.south)+(0.08,-0.12)$) {trap\_fin\_i};
\node[puerta y=nnn, y3, anchor=input 2] (gme) at (5.0,2.57) {};
\draw[flujo] (mi.east) -| ($(gme.input 1)+(-0.25,0)$) -- (gme.input 1);
\draw[flujo] (mf.east) -| ($(gme.input 3)+(-0.25,0)$) -- (gme.input 3);
\red[0.6]{gme.input 2}{acumula}
\redsal{gme.output}{en\_meseta}
\node[op] (ma2) at (9.0,2.95) {$+$};
\red{ma2.west}{trap\_i}
\redab[0.2]{ma2.south}{}
\node[red, anchor=west] at ($(ma2.south)+(0.08,-0.12)$) {suma};
\node[op] (mc2) at (11.2,2.2) {$+1$};
\red{mc2.west}{cuenta}
\node[reg=1.1cm] (rsu) at (13.6,2.95) {suma};
\node[reg=1.1cm] (rcu) at (13.6,2.2) {cuenta};
\draw[flujo] (ma2.east) -- (rsu.west);
\draw[flujo] (mc2.east) -- (rcu.west);
\node[chico, anchor=south] at ($(rsu.north)+(0,0.03)$) {en: \texttt{en\_meseta}; 0 con \texttt{abre}};
\sal{rsu.east}{trap\_sum\_o}
\sal{rcu.east}{trap\_n\_o}

% ========================================================= marcas de solape ===
% durante la medición: otro disparo adentro, el primero fija dt_next
\node[reg=1.1cm] (rad) at (5.0,0.5) {adentro};
\node[reg=1.1cm] (rdn) at (5.0,-0.65) {dt\_next};
\node[reg=1.1cm] (rpe) at (5.0,-1.8) {perdido};
\red{rad.west}{1}
\red{rdn.west}{len}
\red{rpe.west}{1}
\node[chico, anchor=south] at ($(rad.north)+(0,0.03)$) {en: \texttt{acumula} $\wedge$ \texttt{disparo\_i}; 0 con \texttt{abre}};
\node[chico, anchor=south] at ($(rdn.north)+(0,0.03)$) {en: \texttt{acumula} $\wedge$ \texttt{disparo\_i} $\wedge$ \texttt{dt\_next} $=$ máx; máx con \texttt{abre}};
\node[chico, anchor=south] at ($(rpe.north)+(0,0.03)$) {en: \texttt{acumula} $\wedge$ \texttt{perdido\_i}; 0 con \texttt{abre}};
% al abrir: el contexto del disparo, que dejó el asignador
\node[reg=1.1cm] (rwi) at (11.0,-0.05) {win\_id};
\node[reg=1.1cm] (rdp) at (11.0,-1.2) {dt\_prev};
\node[reg=1.1cm] (rot) at (11.0,-2.35) {otras};
\red{rwi.west}{win\_id\_i}
\red{rdp.west}{dt\_prev\_i}
\red{rot.west}{otras\_i}
\node[chico, anchor=south] at ($(rwi.north)+(0,0.03)$) {se cargan con \texttt{abre}};
\sal{rad.east}{adentro\_o}
\sal{rdn.east}{dt\_next\_o}
\sal{rpe.east}{perdido\_o}
\sal{rwi.east}{win\_id\_o}
\sal{rdp.east}{dt\_prev\_o}
\sal{rot.east}{otras\_o}
% solapada: alguna abierta al abrir, o un disparo adentro
\node[op] (nzo) at (14.6,-2.85) {$\neq0$};
\red{nzo.west}{otras}
\node[puerta o, y2, anchor=input 2] (gso) at ($(nzo.east)+(0.3,0)$) {};
\draw[flujo] (nzo.east) -- (gso.input 2);
\red[0.9]{gso.input 1}{adentro}
\sal{gso.output}{solapada\_o}

% ============================================================== el límite ===
\node[contenedor, fit={(0.15,14.9) (19.3,-3.55)}] (lim) {};
\node[titulo, anchor=south west, inner sep=2pt, align=left] at (lim.north west)
  {\texttt{\bfseries mca\_slot}\;\; {\scriptsize\itshape expandido}\\
   \prm{DW=14 & QW=32 & AMP\_W=16 & LEN\_W=16\\ PUB\_IDW=16 & NW=4 & TW=18 & TNW=8\\ NF=3}};

% ================================================================= leyenda ===
\leyenda{0.2,-4.3}{amarillo!22/segmentación, verde!14/estadísticos,
  grisr!45/contabilidad, grisr/registro}{}

\end{tikzpicture}
\end{document}
